\begin{abstract}
Synthetic records can be individually plausible while collectively forming a
poor dataset: accepted records may repeat one another, concentrate around a
small region of support, omit important subgroups, or fail to improve a
downstream task. We formulate \rcqasyn{} as a research direction that treats
synthetic-data creation as sequential dataset construction rather than
one-shot generation followed by a single filter. Its core separates
record-level validity from dataset-level quality, represents validator outputs
as scoped heterogeneous evidence, and proposes calibrated fusion for
risk-aware acceptance decisions. We describe a reproducible evaluation design
spanning deterministic rules, an uncalibrated language-model judge, copying
and support diagnostics, human record-validity annotation, intrinsic metrics,
and downstream utility. This provisional manuscript reports the implemented
experimental foundation and development observations; human annotation is in
progress, and calibration, risk-aware control, and final benchmark evaluation
remain pending. Localized repair and richer dataset optimization are possible
later extensions rather than experimentally validated components.
% TODO[RESULTS]: replace after final benchmark and completed human evaluation.
\end{abstract}
